\documentclass{article}
\usepackage{amsmath,amssymb,mathrsfs}
\usepackage[margin=1in]{geometry}
\allowdisplaybreaks
\begin{document}

\section*{CSS-based distributed TMPC: local problem $\mathscr P_i(t)$}
For each following CAV $i=1,\ldots,N$, the local predictive-control problem is
\begin{equation}
\min_{\bar u_i^p(k\mid t),\,k\in\mathbb K_{0:N_p-1}}
J_i\bigl(\bar x_i^p,\bar u_i^p,S_i^p,\bar x_i^a;\mathbb K_{0:N_p-1}\mid t\bigr)
=\sum_{k=0}^{N_p-1}l_i\bigl(\bar x_i^p,\bar u_i^p,S_i^p,\bar x_i^a;k\mid t\bigr).
\tag{14a}
\end{equation}
Subject to
\begin{align}
\bar x_i^p(0\mid t)&=\bar x_i(t),\tag{14b}\\
\bar x_i^p(k+1\mid t)&=A_d\bar x_i^p(k\mid t)+B_d\bar u_i^p(k\mid t),\tag{14c}\\
\bar e_i^p(k\mid t)&=\frac{1}{|\mathbb I_i|}\sum_{j\in\mathbb I_i}
\left(C\bar x_i^p(k\mid t)-C\bar x_j^a(k\mid t)
 +(i-j)\left(l+h_iD\bar x_i^p(k\mid t)\right)\right),\tag{14d}\\
s_i^p(k\mid t)&=\bar e_i^p(k\mid t)
 +\sum_{\tau=0}^{t-1}k_i\bar e_i(0\mid\tau)\Delta t
 +\sum_{\tau=0}^{k}k_i\bar e_i^p(\tau\mid t)\Delta t,\tag{14e}\\
S_i^p(k\mid t)&=\beta^{i-f}s_f^a(k\mid t)-s_i^p(k\mid t),\tag{14f}\\
\bar u_i^p(k\mid t)&\in\overline{\mathbb U}_i,\tag{14g}\\
\bar x_i^p(k\mid t)&\in\overline{\mathbb X}_i,\tag{14h}\\
S_i^p(N_p\mid t)&=0.\tag{14i}
\end{align}
The source prints these rows without separate row-wise quantifiers; its trajectory domains below give $k\in\mathbb K_{0:N_p-1}$ for inputs and the dynamics transitions, and $k\in\mathbb K_{0:N_p}$ for states, errors, and sliding variables.

\subsection*{Stage cost, notation, and trajectory roles (physical pages 3 and 7)}
\begin{equation}
\begin{split}
l_i\bigl(\bar x_i^p,\bar u_i^p,S_i^p,\bar x_i^a;k\mid t\bigr)
={}&\|S_i^p(k\mid t)\|_{Q_i}
 +\|\bar u_i^p(k\mid t)\|_{R_i}\\
&+\|\bar x_i^p(k\mid t)-\bar x_i^a(k\mid t)\|_{F_i}.
\end{split}\tag{15}
\end{equation}
Here $Q_i,R_i,F_i$ are positive weight matrices. The notation section defines
\[
\|x\|_Q\triangleq\sqrt{x^{\mathsf T}Qx},
\qquad x\in\mathbb R^n,\quad Q\in\mathbb R^{n\times n}
\text{ positive definite}.
\]
Thus (15) contains weighted norms, not squared weighted norms. Also,
\[
C=\begin{bmatrix}1&0\end{bmatrix},\qquad
D=\begin{bmatrix}0&1\end{bmatrix}.
\]
For arbitrary integers $N_1<N_2$, the source defines
\[
\mathbb K_{N_1:N_2}=\{N_1,N_1+1,\ldots,N_2-1,N_2\}.
\]
The notation $x(k\mid t)$ denotes the value of $x(t+k)$ predicted at time $t$. The predictive horizon is $N_p$; within $[t,t+N_p]$ the source distinguishes
\[
\begin{array}{lll}
\text{predicted:}&
\bar x_i^p(k\mid t),\bar e_i^p(k\mid t),s_i^p(k\mid t),S_i^p(k\mid t),
&k\in\mathbb K_{0:N_p},\\
&\bar u_i^p(k\mid t),&k\in\mathbb K_{0:N_p-1},\\[2pt]
\text{optimal:}&
\bar x_i^*(k\mid t),\bar e_i^*(k\mid t),s_i^*(k\mid t),S_i^*(k\mid t),
&k\in\mathbb K_{0:N_p},\\
&\bar u_i^*(k\mid t),&k\in\mathbb K_{0:N_p-1},\\[2pt]
\text{assumed:}&
\bar x_i^a(k\mid t),\bar e_i^a(k\mid t),s_i^a(k\mid t),
&k\in\mathbb K_{0:N_p},\\
&\bar u_i^a(k\mid t),&k\in\mathbb K_{0:N_p-1}.
\end{array}
\]
The predicted quantities parameterize the optimization; the starred quantities are its solution. The assumed trajectories are generated from optimal trajectories and communicated to out-neighbors. In a local solve, the current nominal state, error history in (14e), assumed trajectories, topology, and controller parameters are data; the minimized-over variables are the predicted nominal inputs. Predicted states, errors, and sliding variables are linked to those inputs by (14b)--(14f).

\subsection*{Topology and physical inputs (physical pages 3--4)}
There are $N$ following CAVs and a virtual leader $0$, on a flat, straight road. Communication delays, packet losses, and failures are neglected. The following-vehicle graph and augmented graph are
\[
\mathscr G=(\mathscr V,\mathscr E),\qquad
\mathscr V=\{1,\ldots,N\},\qquad
\overline{\mathscr G}=(\overline{\mathscr V},\overline{\mathscr E}).
\]
The augmented graph includes the virtual leader. The set $\mathbb I_i$ consists of all neighbors from which CAV $i$ can receive information, and $\mathbb O_i$ consists of all neighbors to which it can send information. The symbol $|\mathbb I_i|$ is the cardinality of the in-neighbor set.

\paragraph{Assumption 1.}
The augmented graph is a directed acyclic graph (a finite directed graph without directed cycles) and contains a directed spanning tree rooted at node $0$. The source's accompanying ``i.e.'' clause is retained as printed:
\[
n_0=0,n_1,\ldots,n_N=i
\quad\text{is a sequence of all the nodes such that}\quad
(n_i,n_{i+1})\in\overline{\mathscr E},
\quad\forall i\in\{0,1,2,\ldots,N-1\}.
\]
\emph{Printed scope issue:} the clause reuses $i$ as endpoint and running index, and a sequence through all nodes is not generally equivalent to existence of a rooted spanning tree. Neither the clause nor the stated DAG/rooted-tree premise is replaced here.

For $i=0,1,2,\ldots,N$, the actual vehicle dynamics are
\begin{equation}
\dot p_i(t)=v_i(t),\qquad
\dot v_i(t)=u_i(t)+d_i(t).
\tag{1}
\end{equation}
Here $p_i,v_i,u_i,d_i$ are position, velocity, desired acceleration, and lumped unknown disturbance, respectively. The disturbance is unknown but bounded within a compact interval:
\[
d_i(t)\in\mathbb W_i,\qquad \mathbb W_i\text{ is a compact interval}.
\]
Without disturbance, the nominal quantities satisfy
\begin{equation}
\dot{\bar p}_i(t)=\bar v_i(t),\qquad
\dot{\bar v}_i(t)=\bar u_i(t).
\tag{2}
\end{equation}
The actual and nominal state vectors and their deviation are
\[
x_i(t)=\begin{bmatrix}p_i(t)&v_i(t)\end{bmatrix}^{\mathsf T},\qquad
\bar x_i(t)=\begin{bmatrix}\bar p_i(t)&\bar v_i(t)\end{bmatrix}^{\mathsf T},\qquad
\widetilde x_i(t)=x_i(t)-\bar x_i(t).
\]
\begin{equation}
\dot{\bar x}_i(t)=A\bar x_i(t)+B\bar u_i(t),\qquad
A=\begin{bmatrix}0&1\\0&0\end{bmatrix},\quad
B=\begin{bmatrix}0\\1\end{bmatrix}.
\tag{3}
\end{equation}
Euler discretization with sampling time $\Delta t$ gives
\begin{equation}
\bar x_i(k+1)=A_d\bar x_i(k)+B_d\bar u_i(k),\qquad
A_d=\begin{bmatrix}1&\Delta t\\0&1\end{bmatrix},\quad
B_d=\begin{bmatrix}0\\\Delta t\end{bmatrix}.
\tag{13}
\end{equation}

\subsection*{Spacing errors and coupled sliding surface (physical pages 4 and 6)}
The fixed spacing-policy parameters $l$ and $h_i$ are the predefined constant distance and time headway, respectively. The actual and nominal average spacing errors are
\begin{align}
e_i(t)&=\frac{1}{|\mathbb I_i|}\sum_{j\in\mathbb I_i}
\left(p_i(t)-p_j(t)+(i-j)(l+h_iv_i(t))\right),\tag{4}\\
\bar e_i(t)&=\frac{1}{|\mathbb I_i|}\sum_{j\in\mathbb I_i}
\left(\bar p_i(t)-\bar p_j(t)+(i-j)(l+h_i\bar v_i(t))\right).\tag{5}
\end{align}
The accompanying velocity errors are
\[
e_{vi}(t)=v_i(t)-v_0(t),\qquad
\bar e_{vi}(t)=\bar v_i(t)-\bar v_0(t).
\]
The sliding variables are defined by
\begin{align}
s_i(t)&=\bar e_i(t)+\int_0^t k_i\bar e_i(\tau)\,d\tau,
\qquad k_i>0,\tag{8}\\
S_i(t)&=\beta^{i-f}s_f(t)-s_i(t),\qquad 0<\beta<1.\tag{9}
\end{align}
Specifically,
\[
s_0(t)=0.
\]
The parameter $\beta$ is constant, and $f$ is the index of the nearest in-neighbor CAV. On the string-stable CSS, the source states
\[
S_i(t)=0,\qquad s_i(t)=0,
\]
and obtains
\begin{equation}
\dot{\bar e}_i(t)=-k_i\bar e_i(t).
\tag{10}
\end{equation}
The corresponding equivalent control input is
\begin{equation}
u_{si}(t)=
\frac{\displaystyle\sum_{j\in\mathbb I_i}
\left(-\bar v_i(t)+\bar v_j(t)\right)-|\mathbb I_i|k_i\bar e_i(t)}
{\displaystyle h_i\sum_{j\in\mathbb I_i}(i-j)}.
\tag{11}
\end{equation}
These continuous-time relations define the CSS used in the terminal construction; they are not additional stage constraints requiring the predicted trajectory to lie on the CSS before the terminal point.

\subsection*{Feedback and tube-tightened domains (physical pages 6--7)}
The integrated actual control input is
\begin{equation}
u_i(t)=\bar u_i(t)+k_{xi}\widetilde x_i(t),
\qquad k_{xi}\in\mathbb R^{1\times2}.
\tag{12}
\end{equation}
The nominal input $\bar u_i(t)$ is the feedforward term produced by the nominal DMPC. The feedback gain must be selected to guarantee stability of $\widetilde x_i(t)$. The feedback term is intended to keep the actual state within a robustly positively invariant tube, denoted by $\mathbb Z_{xi}$ and centered at the nominal state. The source does not supply a defining invariance equation or a numerical construction for this set.

The tightened domains printed below (14) are
\[
\boxed{\quad
\overline{\mathbb U}_i
=\mathbb U_i\ominus k_{xi}\mathbb Z_{xi}\ominus k_{di}\mathbb W_i,
\qquad
\overline{\mathbb X}_i=\mathbb X_i\ominus\mathbb Z_{xi}.
\quad}
\]
Here $\mathbb U_i$ and $\mathbb X_i$ are the actual control-input and state constraint sets. They remain abstract. The operation $\ominus$ denotes Pontryagin difference; $\oplus$ denotes Minkowski sum.

\emph{Printed omission:} the extra coefficient $k_{di}$ in the input tightening has no definition in the supplied paper, and (12) contains no corresponding disturbance-feedback term. The printed tightening is retained without deleting $k_{di}\mathbb W_i$, assigning a value to $k_{di}$, or modifying (12). The page-7 image reads $k_{xi}$ in the first subtraction, consistently with (12), rather than the $k_{ci}$ given in the auxiliary image transcription.

\subsection*{Assumed trajectories and terminal extension (physical pages 7--9)}
The selected terminal-state extension is the displayed equation (16):
\begin{equation}
\bar x_i^a(N_p\mid t+1)
=A_d\bar x_i^a(N_p\mid t)+B_d\bar u_i^*(N_p\mid t).
\tag{16}
\end{equation}
Its terminal virtual input, defined from (11), is
\begin{equation}
\bar u_i^*(N_p\mid t)=
\frac{1}{\displaystyle h_i\sum_{j\in\mathbb I_i}(i-j)}
\left(
\sum_{j\in\mathbb I_i}
\left(-D\bar x_i^a(N_p\mid t)+D\bar x_j^a(N_p\mid t)\right)
-|\mathbb I_i|k_i\bar e_i^a(N_p\mid t)
\right).
\tag{17}
\end{equation}
This is a virtual input characterizing the terminal assumed state, not an additional optimization variable at $k=N_p$. The terminal assumed spacing error is defined on physical page 9 by
\[
\bar e_i^a(N_p\mid t)=\frac{1}{|\mathbb I_i|}\sum_{j\in\mathbb I_i}
\left(C\bar x_i^a(N_p\mid t)-C\bar x_j^a(N_p\mid t)
 +(i-j)\left(l+h_iD\bar x_i^a(N_p\mid t)\right)\right).
\]
The accompanying terminal sliding-variable extension is
\[
s_i^a(N_p\mid t+1)=0.
\]
The source attributes this extension to the flatness property of the sliding surface, together with (14i) and (17).

Algorithm 1 gives the nonterminal shifts and assumed-input update, for all $i=1,\ldots,N$:
\[
\bar u_i^a(k\mid t+1)=
\begin{cases}
\bar u_i^*(k+1\mid t),&k\in\mathbb K_{0:N_p-2},\\[3pt]
\displaystyle
\frac{\sum_{j\in\mathbb I_i}\left(-D\bar x_i^a(N_p\mid t)+D\bar x_j^a(N_p\mid t)\right)
-|\mathbb I_i|k_i\bar e_i^a(N_p\mid t)}
{h_i\sum_{j\in\mathbb I_i}(i-j)},&k=N_p-1,
\end{cases}
\]
\[
\bar x_i^a(k\mid t+1)=\bar x_i^*(k+1\mid t),
\qquad k\in\mathbb K_{0:N_p-1},
\]
\[
s_i^a(k\mid t+1)=
\begin{cases}
s_i^*(k+1\mid t),&k\in\mathbb K_{0:N_p-1},\\
0,&k=N_p.
\end{cases}
\]
Equation (16), rather than the alternative terminal-state row in Algorithm 1, supplies the terminal state for this selected version.

\paragraph{Initialization and application (Algorithm 1, physical page 8).}
For each $i=1,\ldots,N$, initialize
\[
\bar u_i^a(k\mid0)=0,\qquad k\in\mathbb K_{0:N_p-1},
\]
\[
\bar x_i^a(k\mid0)=
\begin{cases}
x_i(t),&k=0,\\
A_d\bar x_i^a(k-1\mid0)+B_d\bar u_i^a(k-1\mid0),
&k\in\mathbb K_{1:N_p},
\end{cases}
\qquad
\widehat v_i(0)=v_i(0).
\]
At time $t\geq0$, assumed states and sliding variables for $k\in\mathbb K_{0:N_p}$ are transmitted to $j\in\mathbb O_i$ and received from in-neighbors. Each CAV solves $\mathscr P_i(t)$, prepares the shifted assumed trajectories above, and uses
\[
\bar x_i(t+1)=A_d\bar x_i(t)+B_d\bar u_i^*(0\mid t),
\qquad
\widetilde x_i(t)=x_i(t)-\bar x_i(t),
\]
\[
u_i(t)=\bar u_i^*(0\mid t)+k_{xi}\widetilde x_i(t).
\]

\paragraph{Source discrepancies and verification limits.}
Algorithm 1 instead prints the following terminal branch for its assumed-state update:
\[
\bar x_i^a(k\mid t+1)
=A_d\bar x_i^a(N_p-1\mid t+1)
 +B_d\bar u_i^*(N_p-1\mid t+1),\qquad k=N_p.
\]
It uses a next-time optimal input rather than the virtual input of (16)--(17). This conflicting branch is recorded here only as a source note and is not imposed in the selected equation-(16) version. The initialization retains the printed actual state $x_i(t)$ and its time argument; no equality between actual and nominal initial states is supplied. Algorithm 1 initializes $\widehat v_i$ without defining its subsequent role and does not give an explicit initial assumed-sliding-trajectory row. No initialization repair is supplied. Its receive instruction also prints ``from CAVs $i\in\mathbb I_i$'' while displaying the received trajectories with index $j$; the in-neighbor definition supplies the communication role stated above.

The paper supplies neither general numerical bounds for $\mathbb U_i,\mathbb X_i,\mathbb W_i$ nor a special case for a zero denominator in (11) or (17). Its index-set definition is stated only for $N_1<N_2$, although the algorithm uses shortened horizon ranges. No extra bounds, denominator restrictions, or empty-range conventions have been added. The history term $\bar e_i(0\mid\tau)$ in (14e) is retained with its printed lack of a trajectory superscript.

The external detailed neighborhood definitions in Bian et al. (2020a), the feedback/RPI construction references Chen et al. (2022) and Feng et al. (2020), and the cited sliding-surface flatness reference Hansen et al. (2017) are unverified beyond what this paper states. No official supplementary material was supplied. Simulation-specific weights, time steps, horizons, spacings, and acceleration limits are not part of this abstract local formulation.

\end{document}
